\documentclass{article}
\usepackage[T1]{fontenc}
\usepackage{lmodern}
\usepackage{graphicx}
\usepackage{amsmath,amssymb}
\usepackage[a4paper,margin=2cm]{geometry}
\usepackage[absolute,overlay]{textpos}
\setlength{\TPHorizModule}{1mm}
\setlength{\TPVertModule}{1mm}
\usepackage[indonesian]{babel}
\usepackage{hyperref}
\setlength{\emergencystretch}{2em}
% unit: Halaman 13

\begin{document}

% === Halaman 13 (llm) ===

\begin{itemize}
  \item Contoh:
  \begin{itemize}
    \item Plainteks: \texttt{10100010001110101001}
    \item Bagi plainteks menjadi blok-blok 4-bit:
    \texttt{1010 0010 0011 1010 1001}
    ( dalam notasi HEX :A23A9)
  \end{itemize}
  \item Kunci (juga 4-bit): \texttt{1011}
  \item Misalkan fungsi enkripsi $E$ yang sederhana algoritmanya sbb:
  \begin{enumerate}
    \item XOR-kan blok plainteks $P_i$ dengan $K$
    \item geser secara \textit{wrapping} bit-bit dari hasil langkah 1 satu posisi ke kiri
  \end{enumerate}
\end{itemize}

\[ E_K(P) = (P \oplus K) \ll 1 \]

\end{document}
